% 2.1.4
% Ficheiro incluído por main.tex

\subsection[Kerberos: configurar o realm e obter o ticket]{Kerberos: configurar o realm e obter o ticket}

\passo{Passo 1 - Guardar uma cópia do ficheiro original}

\begin{lstlisting}[style=bashstyle]
sudo cp /etc/krb5.conf /etc/krb5.conf.original 2>/dev/null
\end{lstlisting}

\passo{Passo 2 - Escrever o novo \texttt{krb5.conf}}

\begin{lstlisting}[style=bashstyle]
sudo tee /etc/krb5.conf > /dev/null << 'EOF'
[libdefaults]
    default_realm = ISEC.LOCAL
    dns_lookup_realm = false
    dns_lookup_kdc = true
    rdns = false

[realms]
    ISEC.LOCAL = {
        kdc = SR2-SRV.isec.local
        admin_server = SR2-SRV.isec.local
    }

[domain_realm]
    .isec.local = ISEC.LOCAL
    isec.local = ISEC.LOCAL
EOF
cat /etc/krb5.conf
\end{lstlisting}

\emph{Verificação esperada:} o \texttt{cat} mostra o ficheiro exatamente como acima.

\passo{Passo 3 - Obter o ticket}

\begin{lstlisting}[style=bashstyle]
kinit aluno.teste@ISEC.LOCAL
klist
kdestroy
\end{lstlisting}

\emph{Verificação esperada:} o \texttt{kinit} pede a palavra-passe e não mostra mais nada; o \texttt{klist} mostra o utilizador e o ticket \texttt{krbtgt/\allowbreak{}ISEC.LOCAL@ISEC.LOCAL}.

